“Spuit je vrouwtje flink vol
nabij je jagershol!”
“Hou je bek, geile snol.

Ik wil je niet horen.
Je moet me niet storen
bij jou te doorboren.”

Emmy wilde landen
en zeker niet stranden.
“Duw niet met je handen

op mijn neukende kont.”
Emmy legde terstond
haar armen op de grond.

Ze zocht naar oogcontact.
Die waren zonder tact
naar haar poes afgezakt.

Blijkbaar wond haar seksknop
de jager veel meer op.
Dus wachtte ze als Job

op de jagers sperma
en dacht maar aan Dina
van Jakob en Lea.

Dit stuk uit Genesis
leidde tot een crisis
en de besnijdenis

van al Hemor’s mannen.
Het kwam tot die plannen
door het samenspannen

tegen Hemor zijn zoon,
als zijn verdiende loon.
Die vond Dina beeldschoon,

kwam van haar in de ban
en had als dolleman
in het land Kanaän

haar schaamteloos verkracht.
Op de hoogte gebracht
staakte Jacob zijn macht

om wraak te gaan nemen.
Door trouwen te claimen
konden de problemen ook

worden opgelost.
De schuld werd ingelost.
Dina werd uitgedost

voor haar uithuwelijk.
Dat leek wel gruwelijk
als wel afschuwelijk,

maar gewoon voor die tijd.
Het was wel nieuw beleid
dat Jacob had geleid

de trouwrituelen.
“Ieder moest meespelen”,
zat in zijn bevelen.

Hemor’s populatie
kreeg een circumcisie.
Hun recuperatie

bood Jacob de optie
op een repercussie.
Na de executie

namen ze alles mee;
vrouwen, bezit en vee.
Emmy vond het idee

van slinks worden geschaakt,
waarin ze poedelnaakt
haar zuiverheid kwijt raakt,

enigszins bevrijdend.
Ze kwam begeleidend
tot grensoverschrijdend

onzedelijk gedrag.
Zonder het zelfbeklag:
“Dit is niet wat je mag!”

Haar hogere moraal
maakte dan geen kabaal.
Het was niet haar schandaal.

Ze kon afstand nemen
bij kuisheidsproblemen
en haar onschuld claimen.

Ze had zelfs recht op wraak
als vorm van leedvermaak
bij een holle blaaskaak

zoals deze jager.
De gemene plager,
die als een keurslager,

geobsedeerd keek naar
haar bosje rode haar.
Emmy vond het maar raar

dat hij intens genoot
van dat klein stukje bloot
direct boven haar schoot.

Haar vagina jennen
begon al te wennen.
Tijd om te gaan scannen

wat de jager zo dacht.
Was de aantrekkingskracht
dat hij plots had bedacht

zichzelf te behagen?
Ging hij na het jagen
haar om haar hand vragen?

Was haar verovering
zijn sterkste hunkering
in zijn herinnering?

Zag hij door zijn ooglens
haar als een medemens,
een object of een wens?

Zag hij een gestatie,
dus een dracht van Emmy
als een topprestatie?

Voosde hij uit verstand
of stond die aan de kant
en nam hem bij de hand

zijn mannelijk instinct?
Zo werd breed afgevinkt
wat kon worden gelinkt

aan de jagers daden.
Ze moest ernaar raden
om daarmee de kwaden

precies te begrijpen.
Een oog dicht te knijpen
voor alle vergrijpen

van wellicht een vijand.
Dat lag niet voor de hand,
maar zo werd faliekant

elk kwaad overwonnen,
volgens de nonnen
als gedegen bronnen.

Ze keek overwogen
het kwaad in de ogen.
Doch geen mededogen.

“Positiewisseling
is mogelijk het ding
dat zorgt voor inleving.”,

kreeg Emmy daarna door
als vrucht van haar dwaalspoor.
Emmy stelde zich voor

dat ze met hem samen
haar reinheid ontnamen
en diep in haar kwamen.

Maagdelijkheid eren
en deze verteren
behoudens verweren.

Toch kun je, als je schaken als metafoor zou opvatten voor het ontvoeren van een vrouw, enkele parallellen trekken uit het transcript:
Het transcript beschrijft schaken als een spel van manipulatie, macht, strategisch denken, het opzetten van valstrikken en het beheersen of overweldigen van de tegenstander. De speler probeert zijn tegenstander steeds weer te verleiden tot fouten, tot het verliezen van controle over het bord, of zelfs tot psychologisch in het nauw drijven – wat als metafoor gelezen zou kunnen worden als het overrompelen of “overmeesteren” van een ander. Binnen de psychologische dynamiek die wordt aangehaald, zie je elementen als dominantie, manipulatie, het opbouwen van druk, het creëren van afhankelijkheid (door tijdsdruk, bluf en reputatie), en het ontwapenen van verdedigingen.
Al deze strategieën zouden in een bredere, symbolische benadering parallellen kunnen hebben met de tactieken van iemand die een ander tot overgave dwingt, afhankelijk maakt of in de val lokt – wat overeen zou kunnen komen met het archaïsche of literaire gebruik van “schaken” voor het ontvoeren of wegvoeren van een vrouw, bijvoorbeeld zoals de schaakzet “de dame slaan” of “de dame schaakmat zetten”.
Toch blijft het belangrijk te benadrukken dat deze betekenis niet expliciet in het transcript zelf voorkomt. De koppeling tussen “schaken” en het ontvoeren van een vrouw ontstaat hier dus puur op symbolisch, metafoor-niveau, waarbij de machtsstructuren, manipulatie en psychologisch overwicht uit het schaakspel als analogie worden gebruikt voor het onvrijwillig wegnemen of overmeesteren van iemand anders.
Kortom: het transcript geeft je voldoende psychologische strategieën (dominantie, manipulatie, druk, valstrikken, verlies van autonomie) die, wanneer je “schaken” als metafoor voor ontvoering leest, tot een literaire interpretatie kunnen worden omgevormd — maar deze betekenis volgt niet direct uit de tekst zelf en vereist een bewuste verschuiving naar een andere, metaforische lezing.
